\subsection{A parity detector visible in the rank-seventy-three homology cover}

\begin{result}{A.4.14}{Proved}{An explicit homology-visible quaternion parity detector}
Use the literal groups and markings of A.4.12.  In particular
$\Sigma=F(a,b)$, $N_8=\ker(\Sigma\to Q_8)$,
$K_3=\ker(\Sigma\to C_3^2)$, and $T=N_8\cap K_3$ is free
of rank $73$.  Define
\[
 \overline L_8=L_8T',\qquad
 \overline R_8=\overline L_8\cap\Sigma'\Sigma^{12},
 \qquad \overline{\mathcal P}_8=\Sigma/\overline R_8.
                                                        \tag{QH.1}
\]
Here the product is the product of the displayed normal
subgroups.  All these kernels are characteristic in
$\Sigma$, and
\[
 |\Sigma/\overline L_8|=8\cdot3^{24},\qquad
 |\overline{\mathcal P}_8|=32\cdot3^{24}.                \tag{QH.2}
\]
The actual detected modular image $\overline J$, obtained by
the same restriction, deck quotient, and signed abelian
reduction as (QS.8), satisfies
\[
 1\longrightarrow C_3^4\longrightarrow\overline J
 \longrightarrow\operatorname{SL}_2(\mathbb Z/12)
 \longrightarrow1.                                    \tag{QH.3}
\]
Write $\bar c$ for the image of $UL_0$.  Of the $36$
commuting residues modulo twelve, exactly $18$ lift to
$C_{\overline J}(\bar c)$; each surviving fibre has nine
elements.  More precisely, write a commuting residue under
the Chinese remainder isomorphism uniquely as
\[
 q\bmod4=(UL_0)^i\bmod4,\qquad
 q\bmod3=(UL_0)^j\bmod3,\qquad i,j\in\mathbb Z/6.
                                                        \tag{QH.4}
\]
Then its exact lifting condition is
\[
                         i\equiv j\pmod2.              \tag{QH.5}
\]
In particular $5I$, whose two exponents have opposite
parity, does not lift.  The same necessary equality of
parities holds in the actual modular outer surface image
on $\overline{\mathcal P}_8$.

This quotient is an actual quotient of the abelian
rank-$73$ homology cover, joined to the indicated ordinary
abelian target:
\[
 \Sigma/\bigl((T'T^3)\cap\Sigma'\Sigma^{12}\bigr)
              \twoheadrightarrow\overline{\mathcal P}_8.
                                                        \tag{QH.6}
\]
The map is the map induced by the identity of $\Sigma$.
It retains the three-division subgroup $K_3$ and its
quaternion cover $T$ without changing either marking.
\end{result}

\begin{proof}
In the explicit class-two group (QS.6), let
$V_0=N_8/(N_8'N_8^3)\cong\mathbb F_3^9$ denote the
abelian layer.
The two exponent sums modulo three give a surjection
\[
 \ell:V_0\longrightarrow\mathbb F_3^2,
 \qquad W=\ker\ell,\qquad\dim W=7.                     \tag{QH.7}
\]
Surjectivity follows from the exponent lattice
$2\mathbb Z^2$ computed in A.4.12.  Since $L_8\subseteq K_3$,
the image of $T$ in $V=N_8/L_8$ is exactly
$W\oplus\bigwedge^2V_0$.  Its commutator image is exactly
$\bigwedge^2W$: every wedge of two vectors of $W$ is
realized by a commutator of actual lifts in $T$.
This subspace has dimension $\binom72=21$.  Consequently
\[
 N_8/\overline L_8
    =V_0\oplus\bigl(\bigwedge^2V_0/\bigwedge^2W\bigr)
                                                        \tag{QH.8}
\]
with the group law induced from (QS.6).  The two dimensions
are nine and fifteen.  This proves the first order in
(QH.2).  The subgroup $T'$ has zero ordinary exponent sums,
so the exponent lattice of $\overline L_8$ remains
$6\mathbb Z^2$.  The abelian join adds exactly four as in
A.4.12, proving the second order.

Because $T\subseteq N_8$, one has $T^3\subseteq L_8$ and
$T'\subseteq\overline L_8$.  Thus
\[
 (T'T^3)\cap\Sigma'\Sigma^{12}
       \subseteq\overline L_8\cap\Sigma'\Sigma^{12}
       =\overline R_8,                                 \tag{QH.9}
\]
which proves the literal surjective map (QH.6).
The equality
$\overline R_8=R_8T'$ also holds: $T'\subseteq
\Sigma'\Sigma^{12}$, and for an element $lt$ of
$L_8T'$ lying in that abelian kernel the factor $l$ must
lie there too.  Hence $\overline{\mathcal P}_8$ is the
actual quotient of $\mathcal P_8$ obtained by killing
its displayed $T'$ image.  This also fixes its signed
abelianization and its quaternion deck map.

\textbf{The complete finite calculation.}
The program uses the same nine Schreier basis words,
four actual automorphisms, and $9216$ joint linear
representatives as A.4.12.  It computes (QH.7) directly
from their two integer exponent sums modulo three.
The $21$ wedges of a reduced basis of $W$ are independent.
Row reduction supplies the literal projection of the
$36$ central coordinates onto the $15$ coordinates in
(QH.8).  Every one of the four actual automorphisms
preserves the killed subspace, with the matrix identity
checked before passing to the quotient.

The inner maps still have dimension nine.  Thus the
outer central-coordinate space has dimension
$15\cdot9-9=126$.  The retained $36864$ actual Schreier
rows have rank four, with independent zero-based source
indices
\[
                     1356,\ 1384,\ 2000,\ 2028.         \tag{QH.10}
\]
Their span is stable under all four actual generators.
The same complete Schreier argument as in A.4.12 proves
that this is the whole kernel, yielding (QH.3).

For each of the $36$ commuting residues the program
retains its actual positive $u,l$ representative,
its unique quaternion deck commutator, and the exact
version of (QS.15) in the four-dimensional kernel.
Each coefficient matrix has rank two.  Using the six
representatives $q_j$ specified in (QS.16), the complete
cokernel test is
\[
\begin{array}{c|cccccc}
 j&0&1&2&3&4&5\\\hline
 \text{obstruction}&(0,0)&(1,1)&(0,0)&(1,1)&(0,0)&(1,1).
\end{array}                                               \tag{QH.11}
\]
As before, multiplication by an actual power of $\bar c$
bijects the solution fibres in each coset.  All $36$
residues are nevertheless separately evaluated and saved.
There are therefore $18$ surviving residues, each with
$3^{4-2}=9$ lifts, and centralizer order $162$.

For clarity, (QH.4) is also an exact finite computation:
the centralizer of $UL_0$ in each of
$\operatorname{SL}_2(\mathbb Z/4)$ and
$\operatorname{SL}_2(\mathbb Z/3)$ has order six and
consists of its six distinct powers.  These two exhaustive
matrix enumerations are retained in the certificate.
The quotient by simultaneous powers is represented by
$q_j\equiv I\pmod4$, $q_j\equiv(UL_0)^j\pmod3$.
Thus the even classes in (QH.11) are precisely (QH.5).

The producer and complete source certificate are
\path{work/tower_code/modular_quaternion_homology_compressed.py}
and its \path{.json}, \path{.npz}, and \path{.out} files.
In addition to all affine equations and all source words,
the NPZ retains the central projection, the $21$ killed
rows, and the seven-dimensional period kernel.  These
matrices are the comparison maps from the original
quaternion detector; no ambient automorphism or arbitrary
module vector is substituted for the actual modular image.
\end{proof}
